\writeSectionFrame{\presentationTitle}{Singleton im RPG}
\frameWithImageOnly{\BILDERPFAD/rpg1.jpg}
\note{
    In unserem Role Pattern Game benötigen wir eine Klasse, die das Spiel verwaltet. Sie kennt die Spielwelt, die Spieler auf der Welt und den allgemeinen Zustand des Spiels und der Welt. Es bietet sich natürlich an, eine solche Klasse als Singleton zu implementieren, denn wir werden viele Zugriffe darauf benötigen und natürlich benötigen wir immer denselben Spielzustand.
}

\begin{frame}{Anforderung im Rollenspiel}
	\begin{itemize}
 		\item Eine Klasse, die den Zustand des Spiels verwaltet
 		\item Die Klasse muss im gesamten Spiel verfügbar sein
 		\item Es gibt aber nur einen einzigen Spielzustand
 	\end{itemize}
\end{frame}

\begin{frame}[fragile]{Singleton}
	\begin{exampleblock}{\gamepackage{GameManager}}
		\begin{lstlisting}
public class GameManager {
    private static GameManager instance;

    private GameManager() {
        // ...
    }

    public static GameManager getInstance() {
        if (instance == null) {
            instance = new GameManager();
        }
        return instance;
    }
    

    public void startGame() {
    	System.out.println("Spiel startet!");
    }

    public void endGame() {
        System.out.println("Spiel beendet!");
    }
}
	 	\end{lstlisting}
 	\end{exampleblock}
\end{frame}

\begin{frame}[fragile]{Singleton}
	\begin{exampleblock}{\gamepackage{GameManager}}
		\begin{lstlisting}
public class GameManager {
	private static final int MIN_HEIGHT = 10;
	private static final int MIN_WIDTH = 10;
	private static final int MAX_HEIGHT = 30;
	private static final int MAX_WIDTH = 30;
	
    private static GameManager instance;
    private Random random;

    private GameManager() {
        random = new Random();
        
        int width = random.nextInt(MAX_WIDTH);
    	int height = random.nextInt(MAX_HEIGHT);
    	
    	if (height < MIN_HEIGHT) {
    		height = MIN_HEIGHT;
    	}
    	if (width < MIN_WIDTH) {
    		width = MIN_WIDTH;
    	}
    }
}
	 	\end{lstlisting}
 	\end{exampleblock}
\end{frame}
\note{
    Eine Aufgabe dieser Klasse ist es z.B. die Welt zu erzeugen. Damit wir bei jedem Spielstart eine andere Welt bekommen, wird ein einfacher Zufallsgenerierungsmechanismus implementiert. Somit hat der private Konstruktor auch etwas zu tun. 
}